% =============================================================
% Model-necessity study -- contact-aware Jacobian modeling and
% time-variation requirements for the vessel-navigation MPC.
% IEEEtran (journal). Standalone compilable excerpt, companion to
% methodology.tex: the full paper supplies title/abstract/
% introduction/discussion and merges this content where
% appropriate. Written 2026-09-23 from a completed five-part
% offline+Monte Carlo study (Jacobian-mismatch characterization,
% additive-vs-Jacobian residual test, horizon prediction, a
% closed-loop mechanism study that identified and resolved a
% controller-tuning confound, and an n=20 paired Monte Carlo
% robustness campaign).
% =============================================================
\documentclass[journal]{IEEEtran}

\usepackage{amsmath}
\usepackage{amssymb}
\usepackage{amsfonts}
\usepackage{booktabs}

\begin{document}

\section{Model-Necessity Study: Is Contact-Aware, Time-Varying Jacobian Modeling Required?}
\label{sec:model-necessity}

\subsection{Motivation and Hypothesis}
\label{sec:mn-motivation}

The controller of Section~\ref{sec:feedback}\footnote{Cross-reference
into the companion methodology document; section label kept symbolic
here since the two files are not yet merged.} augments a linearized
beam-position prediction with a persistent additive disturbance term
$d_k$ (\S\ref{sec:feedback}). That design implicitly assumes contact's
effect on the model is well approximated by a constant offset. This
section tests that assumption directly.

Let $p_{\mathrm{true}}$ denote the true (contact-aware) tip position and
$p_{\mathrm{model}}$ the internal model's prediction, both linearized
locally around a nominal state and evaluated at a configuration
perturbation $\delta z$:
\begin{equation}
p_{\mathrm{true}} \approx p_0 + J_c\,\delta z,
\qquad
p_{\mathrm{model}} \approx \hat p_0 + J_{nc}\,\delta z,
\end{equation}
where $J_c$ and $J_{nc}$ are the contact-aware and no-contact task
Jacobians evaluated at the same state. The model error decomposes as
\begin{equation}
p_{\mathrm{true}} - p_{\mathrm{model}}
= \underbrace{(p_0 - \hat p_0)}_{\text{additive offset}}
+ \underbrace{(J_c - J_{nc})\,\delta z}_{\text{Jacobian mismatch}}.
\end{equation}
The disturbance term can compensate $d_k \approx p_0 - \hat p_0$, the
first component, but not the second, which depends on the applied
control/configuration change and therefore cannot be captured by a
state-independent additive correction. Two questions follow:
\begin{enumerate}
\item[Q1.] Is contact-induced model error mainly additive, or does
  contact materially rotate the input-output Jacobian?
\item[Q2.] Given the correct contact physics, is a time-varying
  Jacobian required, or is a Jacobian frozen at the trajectory's start
  sufficient?
\end{enumerate}

\subsection{Experimental Design}
\label{sec:mn-design}

A full $2\times2\times2$ factorial isolates these questions. The
nonlinear contact-aware beam model is the truth plant in every
condition; only the controller's internal Jacobian source and
disturbance handling vary:
\begin{center}
\begin{tabular}{@{}lll@{}}
\toprule
Factor & Level 1 & Level 2 \\
\midrule
Jacobian physics & no-contact ($J_{nc}$) & contact-aware ($J_c$) \\
Time dependence & LTI (frozen at start) & LTV (varies along path) \\
Disturbance & $\beta_d=0$ (off) & $\beta_d=1$ (on) \\
\bottomrule
\end{tabular}
\end{center}
giving eight controller-model conditions A--H, with H
($J_c$, LTV, $d$ on) the best-model benchmark and A the deliberately
weakest. All eight share the same task weights, rate/state limits,
disturbance-filter constant, $N=15$, $\Delta t=0.1\,\mathrm{s}$, and
$d=2$-sample delay-aware prediction (joints only; insertion
undelayed, \S\ref{sec:feedback}) established for the validated
controller; only the factorial's three factors are varied.

The reference trajectory is the accepted vessel-navigation plan used
for hardware deployment: a digitized lumen centreline (60 samples,
radius 4.1--5.6\,mm) with the beam base and source magnet raised
$+30\,\mathrm{mm}$ in $z$ and the magnet-exclusion radius set from an
empirically validated closest safe approach, time-parameterized to 180
ticks at $100\,\mathrm{ms}$. Note that the plan's hard contact
constraint (\texttt{contact\_active}) is \emph{never} active at any of
its 118 inverse-configuration nodes: the beam tracks close enough to
the centreline that the binary wall-contact flag does not trigger
anywhere on this trajectory. All contact-related effects reported
below are therefore continuous, sub-contact-threshold model
differences, not consequences of a hard mechanical constraint firing.

\subsection{Jacobian Mismatch Characterization}
\label{sec:mn-exp1}

$J_c(z)$ and $J_{nc}(z)$ were evaluated with matched robot kinematics
along all 118 reference states (finite-difference beam-model
Jacobians in accurate mode; see \S\ref{sec:mn-methods-note} for a fast-mode
artifact caught during this analysis). For the position-only $3\times7$
block, the Frobenius mismatch
$E_J(s)=\lVert J_c-J_{nc}\rVert_F/\lVert J_c\rVert_F$
and the largest principal angle between row spaces,
$\theta_{\mathrm{row}}(s)$, are near zero for the first third of the
path ($s<33\,\mathrm{mm}$: $E_J\approx2\times10^{-5}$,
$\theta_{\mathrm{row}}\approx3\times10^{-4}\,{}^\circ$) and grow
sharply thereafter, reaching $E_J\approx1.9$ and
$\theta_{\mathrm{row}}\approx5.1^\circ$ by $s\approx49\,\mathrm{mm}$.
Row-space rotation stays modest even at peak mismatch (mean
$1.1^\circ$, max $5.1^\circ$): the two models largely agree on
\emph{which} actuator combinations produce task motion. Per-actuator
column rotation is far more informative: joints~1--2 stay well aligned
(mean $3$--$4^\circ$), joints~3--5 rotate substantially (mean
$16$--$28^\circ$, up to $69^\circ$), and the insertion column rotates
the most (mean $20.4^\circ$, \textbf{max $95.1^\circ$} at
$s=49.1\,\mathrm{mm}$) --- past $90^\circ$, contact does not merely
weaken $\partial p/\partial L$, it inverts its effective direction.
Contact therefore acts as a selective, actuator-dependent rotation of
the control mapping rather than a uniform gain change, concentrated in
roughly the trajectory's final third.

\subsection{Additive vs.\ Jacobian-Corrected Residual Test}
\label{sec:mn-exp2}

At ten representative states spanning the full path, a battery of 54
perturbations (14 canonical $\pm$ actuator directions plus 40 random
directions, joint step $0.01\,\mathrm{rad}$, insertion step
$1\,\mathrm{mm}$) was applied to the true contact plant. For each
perturbation $\delta z^{(m)}$ the residual
$r^{(m)}=\Delta p_{\mathrm{contact}}^{(m)} - J_{nc}\,\delta z^{(m)}$
was regressed (36 train / 18 held-out test split) against an additive
model $r\approx d$ and an additive-plus-Jacobian model
$r\approx d+\Delta J\,\delta z$, and the held-out gain
$G_{\Delta J}=1-E_{d+\Delta J}^2/E_d^2$ computed. In the near-null
region ($s<33\,\mathrm{mm}$) $E_d$ is already at the solver's noise
floor ($0.006$--$0.008\,\mathrm{mm}$) and $G_{\Delta J}$ is negative:
fitting $\Delta J$ overfits noise, correctly, since there is no real
signal there. From $s=36\,\mathrm{mm}$ onward, $E_d$ grows to
$0.17$--$0.81\,\mathrm{mm}$ while $E_{d+\Delta J}$ collapses to
$0.02$--$0.08\,\mathrm{mm}$, giving $G_{\Delta J}=0.87$--$0.998$: the
Jacobian correction explains 87--99.8\% of residual variance beyond
what additive compensation alone captures. The independently fitted
$\Delta J$ matches the analytically computed $(J_c-J_{nc})$ from
\S\ref{sec:mn-exp1} to within $4.5$--$12\%$ relative error at four of
five tested nodes, confirming the fit recovers a real local Jacobian
mismatch rather than a curve-fitting artifact. Once mismatch is
non-negligible, additive disturbance compensation is demonstrably
insufficient, and the achievable residual with $d$ alone
($0.17$--$0.81\,\mathrm{mm}$) is on the order of this system's
whole-path tracking targets.

\subsection{Horizon Prediction}
\label{sec:mn-exp3}

All eight conditions were evaluated as pure prediction (linear
integration of the true commanded trajectory; no closed-loop QP)
against the realized nonlinear plant, at horizons
$j\in\{1,3,5,10,15\}$ from eleven starting ticks:

\begin{center}
\begin{tabular}{@{}lrrrrr@{}}
\toprule
Condition & $j{=}1$ & $j{=}3$ & $j{=}5$ & $j{=}10$ & $j{=}15$ (mm) \\
\midrule
A: $nc$, LTI, off & 2.950 & 2.947 & 2.951 & 2.965 & 3.050 \\
B: $nc$, LTI, $+d$ & 0.040 & 0.123 & 0.235 & 0.568 & 0.863 \\
C: $nc$, LTV, off & 2.950 & 2.958 & 2.984 & 2.993 & 3.119 \\
D: $nc$, LTV, $+d$ & 0.040 & 0.122 & 0.230 & 0.497 & 0.736 \\
E: $c$, LTI, off & 0.003 & 0.025 & 0.088 & 0.286 & 0.397 \\
F: $c$, LTI, $+d$ & 0.003 & 0.025 & 0.088 & 0.286 & 0.397 \\
G: $c$, LTV, off & 0.003 & 0.011 & 0.036 & 0.078 & 0.082 \\
H: $c$, LTV, $+d$ & 0.003 & 0.011 & 0.036 & 0.078 & 0.082 \\
\bottomrule
\end{tabular}
\end{center}

($E{\equiv}F$ and $G{\equiv}H$ exactly, by construction, since the
contact model \emph{is} the truth plant and its own disturbance is
zero -- a built-in consistency check.) B illustrates the predicted
textbook signature cleanly: disturbance compensation nearly eliminates
error at $j=1$ ($2.95\to0.04\,\mathrm{mm}$) but it regrows
$\sim\!22\times$ by $j=15$ ($0.04\to0.86\,\mathrm{mm}$) --- the offset
is corrected, the wrong sensitivity accumulates. Paired contrasts at
$j=15$: $H$ vs.\ $D$ ($0.082$ vs.\ $0.736\,\mathrm{mm}$) shows the gap
persists even when both sides have disturbance compensation; $E$ vs.\
$G$ ($0.397$ vs.\ $0.082\,\mathrm{mm}$, $4.9\times$) shows correct
physics still benefits substantially from time-variation, while $A$
vs.\ $C$ (statistically indistinguishable, $3.05$ vs.\
$3.12\,\mathrm{mm}$) shows time-variation provides no benefit at all
when the underlying physics is wrong -- the constant-offset error
dominates regardless. Time-variation and correct contact physics
therefore address distinct, non-interchangeable sources of error.

\subsubsection{Methodological note: a fast-mode Jacobian artifact}
\label{sec:mn-methods-note}

An initial pass of this analysis used the beam model's
\texttt{jacobian\_mode="fast"} finite-difference estimator and showed
an unexplained spike at one specific tick (largest singular value
$\approx79$ versus $\approx1.0$--$1.1$ at every neighboring tick),
while the true tip-position output was smooth there
($0.19\,\mathrm{mm}$ per-tick change, unremarkable). Direct comparison
against \texttt{jacobian\_mode="accurate"} at the same state showed a
smooth singular value ($1.001$), confirming a numerical estimator
artifact rather than real physics. All results in this section use
accurate-mode Jacobians.

\subsection{Closed-Loop Mechanism Study}
\label{sec:mn-exp4}

An initial single-seed closed-loop evaluation of all eight conditions
(genuine receding-horizon QP solves against the nonlinear plant,
command-integrated state propagation) reproduced the horizon-prediction
ranking for six of eight conditions, but showed a reversal: condition D
($J_{nc}$, LTV, $+d$; RMS $0.819\,\mathrm{mm}$) outperformed H
($J_c$, LTV, $+d$; RMS $1.502\,\mathrm{mm}$), and disturbance
compensation \emph{worsened} tracking under LTI (B, F reaching RMS
$\approx5.2\,\mathrm{mm}$, max $\approx8.5\,\mathrm{mm}$, traced to
$74^\circ$ of joint drift by the final tick in a direction the stale
LTI schedule's nullspace projector misclassified as task-irrelevant).
A targeted diagnostic sequence resolved the D-vs-H reversal:

\begin{enumerate}
\item \textbf{Same-state cross-replay.} Re-solving both predictor
  families along each other's realized trajectory showed $J_c$ remains
  the strictly better predictor on \emph{both} D's own and H's own
  states at every horizon (e.g.\ at $j{=}15$: $0.165$ vs.\
  $2.895\,\mathrm{mm}$ on D's trajectory; $0.562$ vs.\
  $2.429\,\mathrm{mm}$ on H's). H's worse closed-loop tracking is
  therefore not attributable to the contact model becoming a worse
  local predictor.
\item \textbf{Exact $Q_N=0$ ablation.} The stagewise nullspace
  posture-regularization term (built from the same schedule Jacobian
  used for prediction) was set to exactly zero (not a scale
  approximation). The reversal disappeared and flipped:
  $D_{\mathrm{noN}}=2.068\,\mathrm{mm}$ vs.\
  $H_{\mathrm{noN}}=0.854\,\mathrm{mm}$, matching the model-fidelity
  ranking from \S\ref{sec:mn-exp1}--\ref{sec:mn-exp3}.
\item \textbf{Projector-source cross experiment.} Decoupling the
  primary prediction Jacobian from the Jacobian used to build the
  nullspace projector (six conditions: $\{J_{nc},J_c\}$ primary
  $\times\{J_{nc},J_c,\text{none}\}$ projector) showed performance
  tracks the \emph{primary} Jacobian ($nc$-primary $\approx0.79$--$0.82\,\mathrm{mm}$;
  $c$-primary $\approx1.49$--$1.50\,\mathrm{mm}$, regardless of
  projector source, which shifts results by only 1--3\%). The
  nullspace-projector \emph{source} is not the mechanism; what matters
  is whether the $Q_N$ term is present at all, and it helps
  $nc$-primary while hurting $c$-primary uniformly. This is most
  consistent with an effective-gain conflict: the more accurate
  contact model's primary tracking correction is more aggressive under
  the fixed $Q_p/R$ weighting, so the fixed-weight posture regularizer
  fights it harder than it fights the gentler no-contact correction.
\item \textbf{$R$-sensitivity check.} With $Q_N$ exactly disabled, a
  closed-loop sweep over $R\in\{100,300,700\}\times R_0$ for the two
  LTV$+d$ conditions confirmed contact wins at every tested $R$ (ratio
  $c/nc$: $0.82$, $0.38$, $0.41$), with the gap stable from
  $R{=}300$ onward, matching the validated $R{=}700R_0$. $R700$ is not
  materially confounding the physics comparison and was frozen for the
  main study.
\end{enumerate}

\subsection{Monte Carlo Robustness Campaign}
\label{sec:mn-mc}

With $Q_N=0$ and $R=700R_0$ frozen (identical $Q_p$, $R_d$, $d=2$,
$\beta_d=1$, $N=15$, $V_f=0$, constraints, disturbance estimator,
reference trajectory, and solver settings across all conditions, no
per-condition retuning), the four core conditions ($J_{nc}^{\mathrm{LTI}}{+}d$,
$J_{nc}^{\mathrm{LTV}}{+}d$, $J_c^{\mathrm{LTI}}{+}d$,
$J_c^{\mathrm{LTV}}{+}d$) were evaluated under matched-seed initial-state
perturbations (joint offsets $\sim\mathcal N(0,2^\circ{}^2)$, insertion
offset $\sim\mathcal N(0,(1\,\mathrm{mm})^2)$, identical realization
across all four conditions per seed, giving a paired rather than
unrelated-groups design). This isolates uncertainty in the initial tip
placement specifically; measurement noise and beam/contact-parameter
uncertainty are not included and are noted as a separate future
campaign rather than mixed into one distribution.

\emph{Time-variation contrasts (10 seeds).} With correct contact
physics, LTV improved RMS by $4.310\,\mathrm{mm}$ on average relative
to LTI (95\% CI $[4.147,4.473]\,\mathrm{mm}$, $10/10$ seeds); with
no-contact physics, by $3.742\,\mathrm{mm}$ (95\% CI
$[3.329,4.157]\,\mathrm{mm}$, $10/10$ seeds). Both effects are large
and unambiguous; ten seeds were sufficient and this contrast was not
extended.

\emph{Contact-modeling contrast (extended to 20 seeds).} The
contact-aware LTV model reduced whole-path RMS tracking error by
$0.666\,\mathrm{mm}$ on average relative to the no-contact LTV model
(95\% CI $[0.430,0.902]\,\mathrm{mm}$), with $16/20$ paired trials
favoring the contact-aware model (median $0.808\,\mathrm{mm}$, IQR
$[0.340,1.097]\,\mathrm{mm}$; standardized paired effect
$d_z=1.32$). The improvement was strongly path-progress dependent.
Splitting the path into three progress-based regimes (early,
$s<33\,\mathrm{mm}$; transition, $33\le s<45\,\mathrm{mm}$; late,
$s\ge45\,\mathrm{mm}$ -- chosen from the empirical $E_J$ transition
in \S\ref{sec:mn-exp1}, \emph{not} from the hard contact-active flag,
which never triggers on this plan, \S\ref{sec:mn-design}):
the early-regime difference was statistically detectable but
practically negligible ($-0.001\,\mathrm{mm}$, 95\% CI
$[-0.0021,-0.0004]\,\mathrm{mm}$, favoring \emph{no-contact} by
roughly one micron); the transition-regime difference was inconclusive
($-0.052\,\mathrm{mm}$, 95\% CI $[-0.146,0.041]\,\mathrm{mm}$); and the
late-regime contact advantage was large and robust
($+1.416\,\mathrm{mm}$, 95\% CI $[0.956,1.877]\,\mathrm{mm}$, $18/20$
trials favoring contact modeling). This directly corroborates
\S\ref{sec:mn-exp1}: the whole-path Monte Carlo advantage is not a
diffuse effect, it is concentrated exactly where the deterministic
Jacobian-mismatch analysis located it.

An exploratory (post-hoc, not used to select seeds) correlation
between the paired difference and the normalized initial-perturbation
magnitude was $\mathrm{corr}(\Delta,\lVert\delta z_0\rVert)=-0.60$:
larger initial departures from the reference are associated with a
smaller or reversed contact-model advantage. A plausible explanation
is that both models' local linearizations degrade together far from
the reference state, but this has not been directly demonstrated and
is reported as a limitation/future-work item, not an established
mechanism.

\subsection{Summary}
\label{sec:mn-summary}

Time-variation is required much more decisively than contact-aware
physics itself: LTV improves tracking by $3.7$--$4.3\,\mathrm{mm}$ in
$10/10$ paired trials regardless of which Jacobian family is used.
Contact-aware modeling provides a smaller but real and robust
additional benefit ($0.67\,\mathrm{mm}$ whole-path, $16/20$ trials)
that is not reproducible by additive disturbance compensation alone
(\S\ref{sec:mn-exp2}--\ref{sec:mn-exp3}) and is concentrated in the
later portion of this vessel's path, where the contact/no-contact
Jacobian mismatch is largest (\S\ref{sec:mn-exp1},
\S\ref{sec:mn-mc}). A controller-tuning confound (a Jacobian-dependent
secondary posture cost) was identified and removed before this
conclusion was drawn (\S\ref{sec:mn-exp4}); the input-reference penalty
weight was checked and found not to be a confound over the tested
range. The perturbation-magnitude sensitivity of the contact-model
advantage is noted as an open question for future work.

\end{document}
